\section{达布--温斯坦定理}
\label{sec:darboux-weinstein}

考虑一般辛流形上的拉格朗日子流形 $L\subset(M,\omega)$。

\begin{definition}\label{def:lagrangian-immersion-embedding}
\stfhSubIndex{I187}{la^ge^lang^ri^zi^liu^xing}{拉格朗日子流形（Lagrangian submanifold）}{qian^ru}{嵌入（embedding）}
\stfhSubIndex{I188}{la^ge^lang^ri^zi^liu^xing}{拉格朗日子流形（Lagrangian submanifold）}{jin^ru}{浸入（immersion）}
设 $(M,\omega)$ 是 $2n$ 维辛流形，$L$ 是维数为
$n=\frac12\dim M$ 的流形。若浸入（相应地，嵌入）$i:L\to M$ 满足
$i^*\omega=0$，则称它为\emph{拉格朗日浸入}（相应地，\emph{拉格朗日嵌入}）。当 $i$ 是嵌入时，称其像为一个（嵌入的）拉格朗日子流形。
\end{definition}

关于拉格朗日子流形，第一个普遍结果是温斯坦 (Wn73) 的如下邻域定理。

\begin{theorem}[达布--温斯坦定理]\label{thm:darboux-weinstein}
\stfhIndexRangeStart{I093}{da^bu^wen^si^tan^ding^li}{达布--温斯坦定理（Darboux–Weinstein theorem）}
设 $j:L\to(M,\omega)$ 是紧拉格朗日浸入，$Z:L\to T^*L$ 是零截面。则存在从 $T^*L$ 中 $Z$ 的一个邻域出发的浸入
\[
 E:V\longrightarrow M
\]
使得
\begin{enumerate}[label=(\arabic*)]
\item $E^*\omega=\omega_0$；
\item $E\circ Z=j$。
\end{enumerate}
若 $j$ 还是嵌入，则可以把 $E$ 取成嵌入。
\end{theorem}

我们把这个定理约化到温斯坦的形变引理，即引理 2.2.22，以此给出证明。这里需要用到\cref{sec:symplectic-linear-algebra}中发展的基本辛线性代数。以下紧随温斯坦 (Wn87) 的叙述。
\stfhIndexRangeEnd{I182}{la^ge^lang^ri^ge^la^si^man^liu^xing}{拉格朗日--格拉斯曼流形（Lagrangian Grassmannian）}

\subsection{达布--温斯坦定理的证明}
\label{subsec:proof-darboux-weinstein}

先给出证明梗概。
\begin{enumerate}[label=(\arabic*)]
\item \textbf{第一步。}使用上一小节发展的典范辛代数，在一次线性化的层面解决寻找微分同胚的问题。更确切地说，我们要在包含映射 $j:L\to M$ 上找到丛映射
$A:T^*L\to TM$，使丛映射
\[
 dj\oplus A:TL\oplus T^*L\longrightarrow j^*TM
\]
成为 $L$ 上的辛丛映射，亦即在各纤维上为辛映射。再把这个辛映射与典范同构
\[
 T(T^*L)|_{o_L}\cong TL\oplus T^*L
\]
复合；将 $o_L$ 与 $L$ 等同后，它便定义恒等映射之上的辛丛映射
\[
 dj\oplus A:T(T^*L)|_{o_L}\longrightarrow j^*(TM).
\]

\item \textbf{第二步。}选择映射 $E_0:V_0\to M$，其中 $V_0$ 是
$o_L\subset T^*L$ 的一个邻域，并使
\[
 dE_0|_{o_L}=dj\oplus A.
\]
把 $V_0$ 取得充分小，$E_0$ 就成为浸入。

\item \textbf{第三步。}比较 $T^*L$ 上的典范辛形式 $\omega_0$ 与
$V_0\subset T^*L$ 上的 $E_0^*\omega$。由 $E_0$ 的选取，
\begin{equation}
 \omega_0|_{o_L}=E_0^*\omega|_{o_L}.
 \label{eq:darboux-weinstein-forms-agree}
\end{equation}
现在应用形变引理 2.2.22，找到微分同胚 $\psi:V\to V_0$，使
$\psi^*(E_0^*\omega)=\omega_0$ 在 $V$ 上成立，并且
$\psi|_{o_L}=\operatorname{id}_{o_L}$。于是浸入 $E=E_0\circ\psi$ 满足要求。
\end{enumerate}

只须完成第一步。注意 $TL\oplus T^*L$ 与 $j^*TM$ 都是 $L$ 上的辛向量丛。我们要用如下引理。

\begin{lemma}\label{lem:lagrangian-splitting-standard-model}
设 $(S,\Omega)$ 是辛向量空间，并设 $S=L\oplus L'$，其中 $L,L'$ 都是拉格朗日的。由
\[
 B(v)=\widetilde\Omega(v)|_L
\]
定义的映射 $B:L'\to L^*$ 是同构，并且直和映射
\[
 1\oplus B:S=L\oplus L'\longrightarrow L\oplus L^*
\]
是辛同构。
\end{lemma}

\begin{proof}
证明留作练习。
\end{proof}

因为 $j:L\to M$ 是拉格朗日浸入，$TL\subset j^*TM$ 是 $L$ 上的拉格朗日子丛。用 $\Omega$ 表示丛 $j^*TM$ 上逐纤维的辛形式。任选光滑补子丛 $U\subset j^*TM$，亦即满足
$TL\oplus U=j^*TM$ 的 $U$。由\cref{lem:lagrangian-complement-average}，
$U^\Omega$ 是另一个光滑补子丛。再取平均
$C=(U+U^\Omega)/2$，它定义一个拉格朗日子丛。这里 $U^\Omega$ 或 $C$ 的光滑性来自第 2.3.8 节中平均的函子性构造。因此我们得到拉格朗日分裂
$j^*TM=TL\oplus C$。把\cref{lem:lagrangian-splitting-standard-model}用于
$C$ 与 $T^*L$，得到丛同构 $B:C\to T^*L$。现令 $A=B^{-1}$，则
\stfhIndexRangeEnd{I093}{da^bu^wen^si^tan^ding^li}{达布--温斯坦定理（Darboux–Weinstein theorem）}
\[
 1\oplus A:TL\oplus T^*L\longrightarrow TL\oplus C=j^*TM
\]
是辛丛同构。这就完成了第一步，进而完成达布--温斯坦定理的证明。

下面陈述达布--温斯坦定理的带参数版本，并把证明留作练习。

\begin{exercise}[带参数的达布--温斯坦定理]\label{ex:parametrized-darboux-weinstein}
\stfhIndex{I223}{can^shu^hua^da^bu^wen^si^tan^ding^li}{参数化达布--温斯坦定理（parameterized Darboux–Weinstein theorem）}
设 $\psi:[0,1]\times L\to M$ 是拉格朗日同痕，即它是光滑映射，且每个
$\psi_t:L\to M$ 都是拉格朗日嵌入。则存在光滑映射
\[
 \Phi:\bigcup_{t\in[0,1]}V_t\subset[0,1]\times T^*L\longrightarrow M
\]
使得 $\Phi_t:V_t\subset T^*L\to M$ 是辛嵌入，并满足
\begin{enumerate}[label=(\arabic*)]
\item $\Phi_t^*\omega=\omega_0$；
\item 对所有 $t\in[0,1]$，$\Phi_t\circ Z=\psi_t$ 在 $L$ 上成立。
\end{enumerate}
特别地，若 $L$ 紧，则可选择邻域
\[
 \mathcal V\subset\overline{\mathcal V}\subset\bigcap_{t\in[0,1]}V_t,
\]
使得对所有 $t\in[0,1]$，$\Phi_t^*\omega=\omega_0$ 在 $\mathcal V$ 上成立。
\end{exercise}
